% !PRINTED 481
\weberpage{481}
%<<CONT
\weq{12}{f(\sqrt{-5})^4 = 1 + \sqrt5,}
\weq{13}{f(\sqrt{-13})^4 = 3 + \sqrt{13}.}
For $m = 17$ the class number is 4, and hence the equation given
above, which is of degree 4 with respect to $u^8$, is the class
equation. If one solves it with respect to $B$, one obtains:
\[
\begin{aligned}
u^8 + \frac{16}{u^8} &= 34 + 10\sqrt{17}, \\
u^4 + \frac{4}{u^4} &= 5 + \sqrt{17}, \\
\left(u^2 + \frac{2}{u^2}\right)^2 &= \frac{(1+\sqrt{17})^2}{2}.
\end{aligned}
\]
If therefore one puts
\[ \sqrt2\,x = f(\sqrt{-17})^2 \]
one obtains for $x$ the quadratic equation:
\weq{14}{x + \frac1x = \frac{1+\sqrt{17}}{2}.}
III. For $\omega$ we put the root of the quadratic
equation:
\[ 2\omega^2 + 2r\omega + n = 0, \]
in which $n$ denotes an odd integer and $r$ an integer
whose square is less than $2n$; thus:
\weq{15}{2\omega + 2r = -\frac{n}{\omega},}
\weq{16}{\omega = \frac{-r+\sqrt{-m}}{2}, \quad m = 2n-r^2.}
Then, by § 34:
\weq{17}{f_2(\omega) f_1(2\omega+2r) = e^{-\frac{r\pi i}{12}}\sqrt2,}
and by (15), (16):
\weq{18}{f_2(\omega) f_2\left(\frac{\omega}{n}\right) = e^{-\frac{r\pi i}{12}}\sqrt2,}
\weq{19}{f_2(\omega) = \frac{e^{-\frac{r\pi i}{12}}\sqrt2}{f_1(r+\sqrt{-m})};}
we therefore put, according as $r$ (and hence $m$) is even or odd,
\weq{20}{x = f_1(\sqrt{-m}), \quad x = f(\sqrt{-m}),}
%CONT>>
